\section{Introduction}
Quadratic unconstrained binary optimization (QUBO) provides a common input for quantum annealers,
quantum-inspired devices and hybrid services. Formulations exist for many combinatorial problems\cite{lucas2014ising}.
However, a person asking where to install charging stations or how to organize an inspection route does not
start with a QUBO matrix. Their request must be clarified, relevant data found and cleaned, a formulation
written, solver resources allocated, and the returned decision explained. Errors at any stage can yield a
feasible solution to a different problem.

Language models can assist this translation. NL4Opt studies natural-language optimization
representations\cite{ramamonjison2023nl4opt}; OptiMUS formulates and solves mixed-integer linear
programs\cite{ahmaditeshnizi2024optimus}. QuantumQUBO Agent decomposes QUBO generation into specialized roles
and evaluates generated formulations on a 100-problem benchmark\cite{mondal2026quantumqubo}.
LLM-QUBO combines language-based formulation with hybrid decomposition\cite{llmqubo2025}.
These lines of work motivate separating modelling from execution and checking. Our focus is the surrounding
workflow: whether the data represents the request, whether the model preserves its requirements, and whether
its output is a usable decision.

The distinction matters in public-data applications. Two tables can use different coordinate systems or spell
the same administrative unit differently. A join by village name can discard entities when names recur across
districts. Separately, a model may enforce a bound of 200 where the user requested five facilities.
Increasing penalties cannot restore missing entities or correct that bound. Consequently, solver feedback must
be interpreted alongside data and formulation evidence before deciding where to repair.

We study the following question: \emph{How can a locally served LLM, supported by specialized roles and executable
feedback, take a confirmed request through data preparation, modelling and annealing while making failures
observable and recoverable?} The operational goal is a feasible, interpretable decision within a bounded budget,
or a concrete explanation of the unresolved gap. A certified optimum is not required for every application.
People confirm the specification, approve data assumptions and review delivery.

Our contributions are threefold:
\begin{enumerate}
  \item A local-LLM workflow with typed data interfaces, tool-supported acquisition and preprocessing,
        versioned artifacts, interchangeable annealing backends and bounded cross-stage repair.
  \item A verification design that separates data correctness, requirement consistency and solver feasibility.
        For small benchmark instances, decoded ground states are evaluated against a checker anchored on the
        benchmark's stated optimum, avoiding dependence on a particular bit layout.
  \item An evidence-separated evaluation: repeated formulation benchmarks, a synthetic end-to-end suite,
        fixed-model backend experiments and specified-source real-data pilots. An implemented autonomous
        acceptance protocol defines the remaining checker, repair and modelling-retrieval comparisons.
\end{enumerate}
The solver-side components reuse established software and encoding methods rather than introduce a new penalty
search algorithm. More agent roles or retrieved patterns are not assumed to help; their incremental effects
require controlled measurements.
